\chapter*{Acrónimos y términos}
\addcontentsline{toc}{chapter}{Acrónimos y términos}

\begin{longtable}{L{0.22\textwidth}L{0.68\textwidth}}
\toprule
\textbf{Sigla o término} & \textbf{Significado} \\
\midrule
\endhead
API & Interfaz de programación de aplicaciones (\emph{Application Programming Interface}). \\
AUC & Área bajo la curva ROC (\emph{Area Under the Curve}). \\
CLI & Interfaz de línea de comandos (\emph{Command-Line Interface}). \\
CSRF & Falsificación de peticiones entre sitios (\emph{Cross-Site Request Forgery}). \\
ENI & Esquema Nacional de Interoperabilidad. \\
ENS & Esquema Nacional de Seguridad. \\
FHIR & \emph{Fast Healthcare Interoperability Resources}, estándar de interoperabilidad de HL7. \\
H & Comparador híbrido con flujo programado. \\
HIS & Sistema de información hospitalaria (\emph{Hospital Information System}). \\
HL7 & \emph{Health Level Seven}, organización que desarrolla estándares de interoperabilidad sanitaria. \\
IA & Inteligencia artificial. \\
JSON & Formato de datos estructurados (\emph{JavaScript Object Notation}). \\
LLM & Modelo de lenguaje de gran escala (\emph{Large Language Model}). \\
M5 & Arquitectura con cinco responsabilidades LLM especializadas. \\
PPI & Inferencia potenciada por predicción (\emph{Prediction-Powered Inference}). \\
R4 & Versión 4 del estándar FHIR. \\
RGPD & Reglamento General de Protección de Datos. \\
S1 & Comparador con un único agente lógico LLM. \\
SAS & Servicio Andaluz de Salud. \\
SMS & Servicio de mensajes cortos (\emph{Short Message Service}). \\
TFM & Trabajo Fin de Máster. \\
\bottomrule
\end{longtable}
